\documentclass[10pt,a4paper]{amsart}
\usepackage[margin=19mm]{geometry}
\usepackage{amsmath,amssymb,mathtools}
\usepackage[scr=rsfs,cal=euler]{mathalfa}
\usepackage{xfrac}
\usepackage[unicode,hidelinks,bookmarks=false]{hyperref}
\newcommand{\ie}{i.e.\ }
\newcommand{\cf}{cf.\ }
\newcommand{\resp}{respectively\ }
\newcommand{\Subsection}{\subsection}
\providecommand{\nobreakdash}{\nobreak\hbox{-}}
\setlength{\emergencystretch}{2em}
\multlinegap0pt
\usepackage{float}
\usepackage{tikz}
\usepackage{amssymb}
\usepackage{mathtools}
\newtheorem{proposition}{Proposition}
\newcommand{\Proj}{\mathrm{P}}
\newcommand{\RePart}{\mathrm{Re}}
\newcommand{\Spec}{\mathrm{Spec}}
\newcommand{\Sph}{\mathbb{S}}
\newcommand{\id}{\mathrm{Id}}
\title{Feed-Forward Steering in Transformer Residual Dynamics}
\author{Timur Mudarisov, Mikhail Burtsev, Radu State}
\date{Source: arXiv:2608.02071v1 (CC BY 4.0)}
\begin{document}
\maketitle

\noindent\emph{Source and licence.} Feed-Forward Steering in Transformer Residual Dynamics. Version arXiv:2608.02071v1 (CC BY 4.0), distributed by arXiv under the Creative Commons Attribution 4.0 International licence, http://creativecommons.org/licenses/by/4.0/, as recorded in the arXiv licence metadata for this exact version and shown on the abstract page https://arxiv.org/abs/2608.02071v1. Full licence notice: \texttt{licenses/arxiv-2608.02071-CC-BY-4.0.txt}.\par\smallskip

\noindent\textit{Editorial scope.} Counted unit: the article's proposition supp:prop:stab (source0.tex lines 1087--1098). The article's complete original proof of it is delivered with it (source0.tex lines 1620--1664, one \texttt{proof} environment of 45 lines, which the article places away from the statement). No mathematics is written, repaired or re-derived: the statement and the proof are the article's own lines, in the article's own notation, and no retained line is substituted or re-flowed.  No further source-local context is needed: every cross-reference of every retained line resolves inside the two delivered spans.  Nothing was omitted from any retained span, so the package carries no omission record.  No retained line contains a citation, so the wrapper carries no bibliography and no bibliography entry.  No retained line contains a figure environment, an image-inclusion command or drawing code, so no image is delivered and the package declares no asset dependency.  Editorial formatting only, in addition: an AMS \texttt{amsart} preamble that restates the article's own declarations and macros used by the retained lines, namely the environments \texttt{proposition} on separate counters, together with the standard packages those lines need.  The retained environments print in this document as Proposition 1 in order of appearance, whereas the article prints the same items with the numbering of its own full document.  The complete native source of the article as distributed in the arXiv e-print for this exact version is bundled unchanged as \texttt{sources/206665/source0.tex}.
\begin{proposition}[Stability test; Proposition 1]\label{supp:prop:stab}
Assume $\Phi$ is $C^1$ in a neighborhood of a critical residual direction
$u_\star$, where $F_\alpha(u_\star)=\rho_\star u_\star$. The intrinsic tangent
Jacobian of $g_\alpha$ at $u_\star$ is
\begin{equation}
J_\star=\Proj_{u_\star}\big(V+\alpha D\Phi(u_\star)-\rho_\star\id\big)
\big|_{T_{u_\star}\Sph^{d-1}}.
\label{supp:eq:jac}
\end{equation}
If $u_\star$ is hyperbolic, then it is locally exponentially asymptotically
stable (a residual attractor) iff $\max\RePart\Spec(J_\star)<0$.
\end{proposition}

\begin{proof}
Let $T_0=T_{u_0}\Sph^{d-1}$ and let
$\psi:T_0\supset\mathcal U\to\Sph^{d-1}$ be the exponential chart centered at
$u_0$, with $\psi(0)=u_0$ and $D\psi(0)=\id_{T_0}$. Define
\begin{equation}
G(v,\alpha)=\Proj_{u_0}g_\alpha(\psi(v))\in T_0.
\label{eq:steering-G}
\end{equation}
At $v=0$, the restriction of $\Proj_{u_0}$ to
$T_{\psi(v)}\Sph^{d-1}$ is the identity. Invertibility is open, so this
restriction remains an isomorphism for all sufficiently small $v$. Hence
$G(v,\alpha)=0$ iff $g_\alpha(\psi(v))=0$ near $(0,0)$.

Because $\Phi$ is $C^2$, $G$ is $C^2$. Proposition~\ref{supp:prop:stab} gives
\begin{equation}
G(0,0)=0,
\quad
D_vG(0,0)=Dg_0(u_0)|_{T_0}=J_0.
\end{equation}
Hyperbolicity implies that $J_0$ is invertible. The $C^2$ implicit function
theorem yields $\alpha_0>0$ and a unique $C^2$ curve $v(\alpha)$ with
$v(0)=0$ and $G(v(\alpha),\alpha)=0$. Put
$u_\alpha=\psi(v(\alpha))$. Differentiating at $\alpha=0$ gives
\begin{equation}
J_0v'(0)+\partial_\alpha G(0,0)=0.
\end{equation}
Since $\partial_\alpha g_\alpha(u)=\Proj_u\Phi(u)$,
\begin{equation}
\begin{aligned}
\partial_\alpha G(0,0)
&=\Proj_{u_0}\Phi(u_0),\\
v'(0)
&=-J_0^{-1}\Proj_{u_0}\Phi(u_0).
\end{aligned}
\end{equation}
Taylor's theorem and $D\psi(0)=\id$ now give
\begin{equation}
u_\alpha
=u_0-\alpha J_0^{-1}\Proj_{u_0}\Phi(u_0)+O(\alpha^2).
\end{equation}
After identifying nearby tangent spaces through the chart, the tangent Jacobian
depends continuously on $\alpha$. Therefore hyperbolicity, and the property that
all eigenvalues have negative real part, persist for sufficiently small
$|\alpha|$.
\end{proof}

\end{document}
